\documentclass[12pt]{article}
\usepackage[margin=1in]{geometry}
\usepackage{enumitem}
\usepackage{titlesec}
\usepackage{parskip}
\usepackage{amsmath}
\usepackage{hyperref}
\usepackage{fontenc}

\title{\textbf{Levy 2006 Claude Guide}\\
\large Levy, Jonah D., ed. \textit{The State after Statism: New State Activities in the Age of Liberalization}.\\
Cambridge, MA: Harvard University Press, 2006.}
\author{}
\date{}

\begin{document}

\maketitle
\tableofcontents
\newpage

%--------------------------------------------------
\section{Central Claims}
%--------------------------------------------------

\begin{itemize}[leftmargin=*, itemsep=4pt]
    \item The volume's central target is the ``retreat of the state'' narrative of neoliberal globalization---the widespread claim that market liberalization, privatization, and global economic integration have simply shrunk the state, reducing its size, reach, and capacity. Levy and his contributors argue this narrative is empirically wrong, or at least badly incomplete.
    \item Instead, advanced-industrial states have undergone \textbf{reconfiguration}: they have shed certain traditional activities (ownership of nationalized industry, detailed dirigiste economic planning, universal and unconditional welfare provision) while simultaneously \textbf{taking on new, often more intensive activities}---new regulatory agencies policing liberalized and privatized markets, targeted social programs aimed at liberalization's losers, active labor-market policies, and market-supporting institutional infrastructure.
    \item Liberalization, on this account, is better understood as a change in the \textbf{mode and target} of state activity---from ownership and direct provision toward regulation, oversight, and targeted compensation---rather than as a simple reduction in state size, spending, or capacity.
    \item The ``state after statism'' is not a smaller state so much as a \textbf{differently organized} state: one that governs markets rather than owning firms, that manages social risk through selective and conditional instruments rather than universal entitlements, and that remains a central, active agent in constructing the very markets that liberalization is often said to have unleashed against it.
    \item This reconfiguration is not merely a defensive reaction to market pressure; it reflects continued political demand for state action---from firms that need functioning regulatory and legal infrastructure to operate in liberalized markets, and from citizens who still expect the state to manage the social dislocations that liberalization produces.
    \item The volume is explicitly comparative and multi-domain: rather than offering a single grand theory, it assembles case studies across financial regulation, welfare-state reform, labor-market policy, and industrial policy to document a recurring pattern of state adaptation that a simple retreat narrative cannot capture.
\end{itemize}

%--------------------------------------------------
\section{Core Theoretical Framework}
%--------------------------------------------------

\begin{itemize}[leftmargin=*, itemsep=4pt]
    \item \textbf{Recasting, not retreating}: The volume's organizing concept is that liberalization prompts states to \emph{recast} their role rather than simply withdraw from it. States exit some functions (industrial ownership, comprehensive planning, universal provision) precisely so they can concentrate capacity on new functions (market regulation, targeted redistribution, oversight of privatized utilities and financial institutions).
    \item \textbf{New regulatory activism}: Privatization and financial liberalization do not eliminate the state's role in markets---they transform it. Where the state once owned firms or banks directly, it now builds and staffs regulatory agencies (competition authorities, financial supervisors, independent regulators of privatized utilities) whose job is to police the newly liberalized market. In many cases the regulatory apparatus created after privatization is more elaborate, more specialized, and more intrusive than the direct managerial oversight it replaced.
    \item \textbf{From universal to targeted social policy}: Rather than dismantling the welfare state wholesale, governments facing fiscal and competitive pressure increasingly rely on targeted, means-tested, or activation-oriented programs (active labor-market policy, targeted tax credits, in-work benefits) aimed specifically at compensating or reskilling the losers from market opening, rather than universal, unconditional entitlements. This is retrenchment in form but not necessarily in the state's overall level of social engagement.
    \item \textbf{Compensation for liberalization's losers}: A recurring mechanism across chapters is that liberalization generates identifiable losers (displaced industrial workers, sectors exposed to new competition), and states respond by constructing new, purpose-built compensatory and active-labor-market instruments---a political logic reminiscent of embedded liberalism, but delivered through more selective, individualized instruments than the mid-century welfare state.
    \item \textbf{Varieties of capitalism and national mediation}: The volume connects directly to the comparative-political-economy literature on how distinct national institutional configurations (coordinated vs.\ liberal market economies, statist vs.\ corporatist traditions) mediate common global pressures in different ways. Liberalization is a shared external pressure, but its domestic translation into policy is filtered through each country's pre-existing institutions, producing reconfiguration rather than convergence on a single liberal model.
    \item \textbf{State capacity as durable, not just persistent}: The theoretical stakes are not merely that ``the state survives''---a weak claim most scholars would accept---but that state \emph{capacity} for governing the economy is actively reconstructed and, in several domains, expanded in scope and technical sophistication even as its instruments change.
\end{itemize}

%--------------------------------------------------
\section{Core Case Studies}
%--------------------------------------------------

\begin{itemize}[leftmargin=*, itemsep=4pt]
    \item \textbf{France}: A paradigmatic dirigiste state that dismantled much of its postwar planning apparatus and nationalized-industry ownership from the 1980s onward, yet built extensive new regulatory institutions for privatized utilities and financial markets, and developed targeted social and labor-market instruments (e.g., RMI-style targeted minimum-income and activation programs) even as universal entitlements were reformed. France illustrates the sharpest contrast between the old statist toolkit and the new regulatory/targeted toolkit.
    \item \textbf{Germany}: Examined through the lens of welfare-state and labor-market reform (the run-up to what would become the Hartz reforms) and the reorganization of coordinated-market-economy institutions (banking supervision, codetermination-adjacent bargaining arrangements) under pressure from European integration and global financial liberalization. Germany shows reconfiguration proceeding more incrementally and consensually, consistent with its coordinated-market-economy institutions.
    \item \textbf{Britain}: The most thoroughly liberalized and privatized of the major European cases, yet also the case where a dense new regulatory state (Ofgem, Ofcom, the Financial Services Authority and its predecessors, and similar sectoral regulators) was constructed essentially from scratch to police the newly privatized utilities and liberalized financial sector---a vivid illustration that privatization begets regulation rather than simply begets absence of the state.
    \item \textbf{Cross-national comparison in social policy}: Chapters on welfare-state adjustment compare how different countries convert a common pressure (fiscal stress, demographic aging, capital mobility) into different mixes of retrenchment, recalibration, and new targeted/activating instruments, tracing the domain-specific and country-specific politics of adjustment rather than assuming a single convergent outcome.
    \item \textbf{Financial regulation and industrial policy chapters}: Additional chapters trace the buildup of financial-market supervision after banking and securities-market liberalization, and the transformation (rather than elimination) of industrial policy from picking-winners subsidization toward competitiveness-oriented, framework-setting, and innovation-support activities.
\end{itemize}

%--------------------------------------------------
\section{Core Works Cited}
%--------------------------------------------------

\begin{itemize}[leftmargin=*, itemsep=4pt]
    \item \textbf{Peter Hall and David Soskice}, \textit{Varieties of Capitalism} (2001)---the dominant comparative-institutionalist framework this volume engages, distinguishing liberal market economies from coordinated market economies and arguing that institutional complementarities channel common pressures (like liberalization) into divergent national outcomes. Levy's volume can be read as extending VoC-style institutionalism specifically to the question of \emph{state} reconfiguration rather than firm/employer coordination.
    \item \textbf{Vivien Schmidt}---work on discourse, ideas, and adjustment in European political economies (e.g., \textit{The Futures of European Capitalism}); Schmidt's emphasis on how political actors discursively legitimate and frame adjustment strategies complements the volume's account of \emph{how} states justify recasting rather than retreating.
    \item \textbf{Suzanne Berger}---work on French industrial and statist traditions, providing essential background on the dirigiste model that French liberalization dismantled and against which the volume's ``recasting'' thesis is measured.
    \item \textbf{John Zysman}---work on French and German finance-led/state-led industrial governance (e.g., \textit{Governments, Markets, and Growth}), a foundational account of how states organize industrial finance that the volume's chapters on financial liberalization directly build on and complicate.
    \item \textbf{Paul Pierson}, \textit{Politics in Time} (2004) and his welfare-state retrenchment scholarship (\textit{Dismantling the Welfare State?}, 1994)---directly relevant to how states adjust rather than simply retreat: Pierson's concepts of \textbf{blame avoidance} and \textbf{path-dependent} welfare-state adjustment explain why governments facing pressure to retrench pursue incremental, targeted, and often obscured reforms rather than wholesale dismantlement---precisely the pattern of ``recasting'' the Levy volume documents in the social-policy domain.
\end{itemize}

%--------------------------------------------------
\section{Methodological Contributions and Limitations}
%--------------------------------------------------

\begin{itemize}[leftmargin=*, itemsep=4pt]
    \item \textbf{Edited-volume, comparative-case methodology}: The book's strength lies in assembling detailed, policy-domain-specific case studies (financial regulation, welfare reform, labor-market policy, industrial policy) across several major European political economies, allowing fine-grained institutional and historical detail that a single-country or single-domain study could not offer.
    \item \textbf{Weaker on a unified causal model}: As with many edited volumes, the tradeoff for this empirical richness is a lack of a single, tightly specified causal model or set of testable propositions that would let one adjudicate \emph{when} reconfiguration takes the form of more regulation versus more targeting versus outright retrenchment. The framework is more a descriptive lens than a predictive theory.
    \item \textbf{The falsifiability problem}: The central conceptual risk is that ``reconfiguration'' (like related ``the state persists but changes form'' arguments) is difficult to falsify. Almost any observed pattern of declining direct state ownership/provision alongside any new state activity---however minor or however clearly compensatory rather than substantive---can be redescribed as evidence of ``reconfiguration'' rather than genuine retreat. Without a clear metric of what would count as disconfirming evidence (i.e., an actual net decline in state capacity), the thesis risks unfalsifiability.
    \item \textbf{Selection on cases with strong state traditions}: The core cases (France, Germany, Britain) are all large advanced-industrial democracies with historically strong administrative states. Whether the reconfiguration pattern generalizes to smaller states, weaker administrative traditions, or developing/middle-income economies undergoing liberalization is left largely unaddressed.
    \item \textbf{Difficulty aggregating across domains}: Because the volume's evidence is domain-specific (finance, welfare, labor markets, industry), it is not obvious how to aggregate these findings into an overall judgment about whether state capacity, on net, expanded, contracted, or merely changed composition---the volume documents change in kind more persuasively than it measures change in degree.
\end{itemize}

%--------------------------------------------------
\section{Connections to the Broader Literature}
%--------------------------------------------------

\begin{itemize}[leftmargin=*, itemsep=4pt]
    \item \textbf{Huntington 1968}: Huntington's concept of \textbf{institutionalization}---the capacity of political institutions to adapt their organization and procedures to new demands over time, rather than remain static or simply collapse under pressure---closely parallels this volume's account of states adapting their forms of activity (new regulatory bodies, new targeted social instruments) rather than simply shrinking under liberalization pressure. Both works treat institutional adaptability, not institutional size, as the key variable.
    \item \textbf{North 1990}: North's \textit{Institutions, Institutional Change and Economic Performance} provides a general theoretical backdrop---path dependence and the incremental, bounded nature of institutional change---for the volume's more specific empirical claims about how state institutions reconfigure (rather than transform wholesale or disappear) under liberalization pressure.
    \item \textbf{Iversen and Soskice 2019}: \textit{Democracy and Prosperity} is a direct thematic sibling. Both works reject simple ``capital disciplines/shrinks the state'' narratives of globalization and instead emphasize that democratic states actively construct the institutional conditions capitalism needs to function (education systems, financial regulation, macroeconomic stabilization). The key difference in emphasis: Iversen and Soskice focus on the advanced knowledge economy and the political coalitions sustaining it over the long run, while Levy's volume focuses more narrowly on the liberalization-era policy-adjustment process itself (the 1980s--2000s) and its domain-specific mechanisms.
    \item \textbf{Boix 2003}: Boix's asset-mobility argument in \textit{Democracy and Redistribution}---that capital mobility disciplines redistributive states by threatening exit---is a useful point of contrast. Where Boix's model implies that mobile capital should press states toward less redistribution and regulation, Levy's volume's institutionally grounded account shows states retaining and even expanding real capacity (regulatory, redistributive) even amid capital mobility and liberalization, suggesting Boix's stylized asset-mobility logic understates the durability of state capacity that comparative-institutionalist accounts emphasize.
    \item \textbf{Ostrom 1990}: Ostrom's account of institutions as durable, adaptable rule-systems that communities design and redesign to manage collective-action problems resonates with the volume's picture of states redesigning their regulatory and social-policy institutions in response to new problems generated by liberalization, rather than institutions simply eroding under external pressure.
    \item \textbf{Putnam 1993}: Putnam's emphasis on the historical persistence and adaptability of institutional capacity (civic community shaping how effectively regional government institutions perform their tasks) offers a loose parallel to this volume's claim that state capacity, once built, can be redirected toward new tasks rather than simply lost; both works push back against purely structural or economistic accounts of institutional change.
    \item \textbf{Przeworski et al.\ 2000}: The \textit{Democracy and Development} project's attention to how political-institutional configurations mediate economic outcomes across many countries parallels this volume's insistence that liberalization's effects on the state are conditioned by national institutions rather than uniform across cases---both resist single deterministic accounts of how markets reshape politics.
    \item \textbf{Muller and Strom 1999}: \textit{Policy, Office, or Votes?} and the broader party-goals/coalition-politics literature supply the missing political mechanism for \emph{how} the welfare-state and labor-market reconfiguration documented in this volume actually gets enacted---party leaders balancing office-seeking, vote-seeking, and policy-seeking goals within governing coalitions determine which specific mix of retrenchment, targeting, and new activism a government pursues, and when. Levy's volume describes the resulting policy patterns; Muller and Strom's framework helps explain the intra-coalition bargaining that produces them.
\end{itemize}

\end{document}
